%%%%%%%%%%%%%%%%%%%%%%%%%%%%%%%%%%%%%%%%%%%%%%%%%%%%%%%%%%%%%%%%%%%%%%
% Lớp 9 | Chương X. Một số hình khối trong thực tiễn | Hình trụ và hình nón
% Dạng 9H0?31-2: Diện tích xung quanh, thể tích hình nón
%%%%%%%%%%%%%%%%%%%%%%%%%%%%%%%%%%%%%%%%%%%%%%%%%%%%%%%%%%%%%%%%%%%%%%


% nguồn: 2026-11-15_09_de_tham_khao_quan_3_de_1_va_de_2 (nhập 15/11/2026)
\begin{bt}%[9H0V31-2]%[ma:714f2580]
% nguon: Đề tham khảo Tuyển sinh 10 TP.HCM 2025-2026 – Quận 3 – Đề 2 – Bài 5
Để làm một chiếc hộp đựng bỏng ngô hình nón, từ một miếng giấy hình tròn bán kính $R = 25\text{ cm}$, bạn An cắt bỏ phần hình quạt tròn có góc ở tâm là $135^\circ$, sau đó dán phần hình quạt lớn còn lại để tạo thành chiếc hộp hình nón.
\begin{enumerate}[a)]
	\item Tính diện tích giấy được sử dụng để làm chiếc hộp trên (làm tròn đến hàng phần trăm của $\text{cm}^2$).
	\item Biết chiếc hộp có thể chứa nhiều hơn thể tích thực của nó tối đa là $15\%$ khi bỏng ngô vun lên phía trên miệng hộp. Tính thể tích bỏng ngô tối đa mà chiếc hộp có thể chứa được (làm tròn kết quả đến hàng đơn vị của $\text{cm}^3$).
\end{enumerate}
\loigiai{
\begin{enumerate}[a)]
	\item Góc ở tâm của phần hình quạt lớn dùng làm hộp nón là:
	$\alpha = 360^\circ - 135^\circ = 225^\circ$.\\
	Diện tích giấy sử dụng để làm hộp chính là diện tích của hình quạt tròn này:
	$S = \dfrac{\pi R^2 \cdot 225^\circ}{360^\circ} = \dfrac{5}{8}\pi \cdot 25^2 = \dfrac{3125\pi}{8} \approx 1227{,}18\text{ (cm}^2\text{)}$.
	\item Độ dài đường sinh của hình nón bằng bán kính của hình quạt tròn: $l = R = 25\text{ cm}$.\\
	Chu vi đáy của hình nón bằng độ dài cung của hình quạt tròn góc $225^\circ$:
	$C = \dfrac{2\pi R \cdot 225^\circ}{360^\circ} = \dfrac{5}{8} \cdot 2\pi \cdot 25 = \dfrac{125\pi}{4}\text{ (cm)}$.\\
	Bán kính đáy $r$ của hình nón thoả mãn: $2\pi r = \dfrac{125\pi}{4} \Rightarrow r = \dfrac{125}{8} = 15{,}625\text{ (cm)}$.\\
	Chiều cao của hình nón là:
	$h = \sqrt{l^2 - r^2} = \sqrt{25^2 - (15{,}625)^2} = \sqrt{625 - 244{,}140625} = \sqrt{380{,}859375} \approx 19{,}5156\text{ (cm)}$.\\
	Thể tích thực của chiếc hộp hình nón là:
	$V_{\text{thực}} = \dfrac{1}{3}\pi r^2 h = \dfrac{1}{3}\pi \cdot (15{,}625)^2 \cdot 19{,}5156 \approx 4989{,}43\text{ (cm}^3\text{)}$.\\
	Thể tích bỏng ngô tối đa chiếc hộp có thể chứa được (nhiều hơn $15\%$ thể tích thực) là:
	$V_{\max} = 115\% \cdot V_{\text{thực}} = 1{,}15 \cdot 4989{,}43 \approx 5737{,}85\text{ (cm}^3\text{)}$.\\
	Làm tròn đến hàng đơn vị, thể tích bỏng ngô chứa được nhiều nhất là khoảng $5738\text{ cm}^3$.
\end{enumerate}
}
\end{bt}

% nguồn: 2026-11-15_09_de_tham_khao_quan_4_de_1_va_de_2 (nhập 15/11/2026)
\begin{bt}%[9H0V31-2]%[ma:da8dfe10]
% nguon: Đề tham khảo Tuyển sinh 10 TP.HCM 2025-2026 – Quận 4 – Đề 2 – Bài 5
Một bể nước rỗng có dạng hình lập phương cạnh $1\text{ m}$. Người ta đặt vào trong khối đó một khối nón đặc bằng bê tông có đỉnh trùng với tâm mặt trên của khối lập phương, đáy hình tròn của khối nón tiếp xúc với bốn cạnh của mặt đáy đối diện. Sau đó người ta đổ nước vào khoảng còn trống giữa hai khối cho đến khi đầy bể. Hỏi cần đổ bao nhiêu lít nước thì đầy khối lập phương? (Lấy $\pi \approx 3{,}14159$, làm tròn kết quả đến hàng đơn vị).
\loigiai{
Thể tích của khối lập phương là:
$V_{\text{hộp}} = 1^3 = 1\text{ (m}^3) = 1000\text{ (lít)}$.\\
Khối nón đặc có:
\begin{itemize}
	\item Chiều cao $h = 1\text{ m}$.
	\item Đáy tiếp xúc với bốn cạnh của hình vuông cạnh $1\text{ m}$ nên đường kính đáy là $1\text{ m}$, suy ra bán kính đáy là $R = 0{,}5\text{ m}$.
\end{itemize}
Thể tích của khối nón đặc là:
$V_{\text{nón}} = \dfrac{1}{3}\pi R^2 h = \dfrac{1}{3} \cdot \pi \cdot (0{,}5)^2 \cdot 1 = \dfrac{\pi}{12}\text{ (m}^3\text{)}$.\\
Thể tích nước cần đổ vào khoảng trống là:
$V_{\text{nước}} = V_{\text{hộp}} - V_{\text{nón}} = 1 - \dfrac{\pi}{12}\text{ (m}^3\text{)}$.\\
Đổi sang đơn vị lít ($1\text{ m}^3 = 1000\text{ lít}$):
$V_{\text{nước}} = 1000 \cdot \left(1 - \dfrac{\pi}{12}\right) \approx 1000 \cdot (1 - 0{,}2618) = 738{,}2\text{ (lít)}$.\\
Làm tròn đến hàng đơn vị, cần đổ khoảng $738\text{ lít}$ nước thì đầy bể.
}
\end{bt}
